% ===============================================================
% AB-06d: Mi ciudad - Wiederholung
% Spanisch Klasse 7 - Sequenz 6: Mi ciudad
% ===============================================================
% Abgeleitet von: LE-06d.tex (Score: 9.2/10)
% Punkte: 20 (4+4+6+6)
% ===============================================================

\documentclass[11pt, a4paper]{article}

% ===============================================================
% SW-MODUS TOGGLE
% ===============================================================
\newif\ifswmodus
\swmodusfalse

% ===============================================================
% PAKETE
% ===============================================================

\usepackage[utf8]{inputenc}
\usepackage[T1]{fontenc}
\usepackage[ngerman]{babel}
\usepackage{helvet}
\renewcommand{\familydefault}{\sfdefault}

\usepackage[margin=2cm]{geometry}
\usepackage{enumitem}
\usepackage{tabularx}
\usepackage{booktabs}
\usepackage{multirow}

\usepackage{xcolor}
\usepackage{framed}
\usepackage{tcolorbox}
\tcbuselibrary{skins}

\usepackage{fancyhdr}
\usepackage{lastpage}
\usepackage{needspace}

% ===============================================================
% FARBEN
% ===============================================================

\definecolor{spanisch-color}{HTML}{c41e3a}
\definecolor{grau-color}{HTML}{888888}
\definecolor{hellgrau-color}{HTML}{f5f5f5}
\definecolor{orange-color}{HTML}{b35c00}

\definecolor{sw-dark}{HTML}{000000}
\definecolor{sw-gray}{HTML}{666666}
\definecolor{sw-white}{HTML}{ffffff}

\ifswmodus
    \colorlet{spanisch}{sw-dark}
    \colorlet{grau}{sw-gray}
    \colorlet{hellgrau}{sw-white}
    \colorlet{orange}{sw-dark}
\else
    \colorlet{spanisch}{spanisch-color}
    \colorlet{grau}{grau-color}
    \colorlet{hellgrau}{hellgrau-color}
    \colorlet{orange}{orange-color}
\fi

\newcommand{\fachfarbe}{spanisch}

% ===============================================================
% VARIABLEN
% ===============================================================

\newcommand{\thema}{Mi ciudad -- Wiederholung}
\newcommand{\sequenz}{06}
\newcommand{\block}{d}
\newcommand{\fach}{Spanisch}
\newcommand{\klasse}{7}
\newcommand{\maxpunkte}{20}
\newcommand{\datum}{\today}

% ===============================================================
% KOPF- UND FUSSZEILE
% ===============================================================

\pagestyle{fancy}
\fancyhf{}
\renewcommand{\headrulewidth}{0.4pt}
\renewcommand{\footrulewidth}{0.4pt}

\lhead{\fach{} -- Klasse \klasse}
\chead{AB-\sequenz\block}
\rhead{\datum}

\lfoot{}
\cfoot{}
\rfoot{Seite \thepage\ von \pageref{LastPage}}

% ===============================================================
% BEFEHLE
% ===============================================================

\newcommand{\punktebox}[1]{\hfill\fbox{\small \_/#1}}
\newcommand{\luecke}[1]{\rule{#1}{0.4pt}}

\newtcolorbox{merkkasten}[1][]{
    colback=hellgrau,
    colframe=\fachfarbe,
    colbacktitle=white,
    coltitle=\fachfarbe,
    boxrule=1pt,
    arc=0pt,
    left=8pt, right=8pt, top=8pt, bottom=8pt,
    fonttitle=\bfseries,
    attach boxed title to top left={yshift=-2mm, xshift=5mm},
    boxed title style={boxrule=0pt, colback=white},
    title=#1
}

\newtcolorbox{vokabelkasten}{
    colback=white,
    colframe=\fachfarbe,
    boxrule=0.5pt,
    arc=0pt,
    left=5pt, right=5pt, top=5pt, bottom=5pt
}

\newtcolorbox{hinweis}{
    colback=white,
    colframe=orange,
    colbacktitle=white,
    coltitle=orange,
    boxrule=1pt,
    arc=0pt,
    left=5pt, right=5pt, top=5pt, bottom=5pt,
    fonttitle=\bfseries,
    attach boxed title to top left={yshift=-2mm, xshift=5mm},
    boxed title style={boxrule=0pt, colback=white},
    title=Hinweis
}

\newenvironment{aufgabe}[1]{%
    \needspace{5\baselineskip}%
    \nopagebreak[4]%
    \vspace{10pt}
    \noindent\textbf{\textcolor{\fachfarbe}{Aufgabe #1}}
    \nopagebreak[4]%
    \vspace{5pt}
    \nopagebreak[4]%
    \begin{list}{}{\leftmargin=0pt}
    \item[]
}{%
    \end{list}
}

\newcommand{\es}[1]{\textcolor{\fachfarbe}{\textbf{#1}}}

% ===============================================================
% DOKUMENT
% ===============================================================

\begin{document}

% ===============================================================
% KOPFBEREICH
% ===============================================================

\begin{center}
    {\Large\bfseries\textcolor{\fachfarbe}{Arbeitsblatt: \thema}}
\end{center}

\vspace{5pt}

\noindent
\begin{tabularx}{\textwidth}{|X|X|X|}
\hline
\textbf{Name:} \luecke{4cm} &
\textbf{Klasse:} \klasse &
\textbf{Datum:} \luecke{2cm} \\
\hline
\end{tabularx}

\vspace{15pt}

% ===============================================================
% MERKKASTEN: Orte in der Stadt
% ===============================================================

\begin{merkkasten}[Merkkasten: Lugares en la ciudad (Orte in der Stadt)]

\begin{tabular}{ll|ll}
\es{la tienda} & = \luecke{3cm} & \es{el supermercado} & = \luecke{3cm} \\[0.5em]
\es{el parque} & = \luecke{3cm} & \es{la farmacia} & = \luecke{3cm} \\[0.5em]
\es{el cine} & = \luecke{3cm} & \es{el hospital} & = \luecke{3cm} \\[0.5em]
\es{la plaza} & = \luecke{3cm} & \es{la estación} & = \luecke{3cm} \\[0.5em]
\end{tabular}

\end{merkkasten}

\vspace{10pt}

% ===============================================================
% MERKKASTEN: Verben ir und estar
% ===============================================================

\begin{merkkasten}[Merkkasten: Die Verben ir und estar]

\begin{tabular}{|l|l|l||l|l|l|}
\hline
\multicolumn{3}{|c||}{\textbf{ir} (gehen, fahren)} & \multicolumn{3}{c|}{\textbf{estar} (sein, sich befinden)} \\
\hline
yo & \luecke{2cm} & voy & yo & \luecke{2cm} & estoy \\
\hline
tú & \luecke{2cm} & vas & tú & \luecke{2cm} & estás \\
\hline
él/ella & \luecke{2cm} & va & él/ella & \luecke{2cm} & está \\
\hline
nosotros & \luecke{2cm} & vamos & nosotros & \luecke{2cm} & estamos \\
\hline
\end{tabular}

\vspace{0.5em}
\textbf{Wichtig:} ir + a + el = ir \es{al} (z.B. Voy \es{al} cine.)

\end{merkkasten}

\vspace{10pt}

% ===============================================================
% MERKKASTEN: Wegbeschreibung
% ===============================================================

\begin{merkkasten}[Merkkasten: Wegbeschreibung]

\begin{tabular}{ll|ll}
\es{todo recto} & = \luecke{3cm} & \es{cerca de} & = \luecke{3cm} \\[0.5em]
\es{a la derecha} & = \luecke{3cm} & \es{lejos de} & = \luecke{3cm} \\[0.5em]
\es{a la izquierda} & = \luecke{3cm} & \es{al lado de} & = \luecke{3cm} \\[0.5em]
\es{Sigue...} & = Geh weiter... & \es{enfrente de} & = \luecke{3cm} \\[0.5em]
\end{tabular}

\end{merkkasten}

\vspace{15pt}

% ===============================================================
% AUFGABE 1: Orte zuordnen
% ===============================================================

\begin{aufgabe}{1 -- Ordne zu: Orte in der Stadt}
Verbinde die spanischen Orte mit der deutschen Bedeutung. \punktebox{4}
\end{aufgabe}

\begin{center}
\begin{tabular}{rcl}
\es{el parque} & \luecke{2cm} & das Kino \\[0.8em]
\es{la tienda} & \luecke{2cm} & der Park \\[0.8em]
\es{el cine} & \luecke{2cm} & die Apotheke \\[0.8em]
\es{la farmacia} & \luecke{2cm} & das Geschäft \\
\end{tabular}
\end{center}

\vspace{10pt}

% ===============================================================
% AUFGABE 2: Verben konjugieren
% ===============================================================

\begin{aufgabe}{2 -- Ergänze die Verben ir und estar}
Schreibe die richtige Form. \punktebox{4}
\end{aufgabe}

\noindent
a) Yo \luecke{2.5cm} (ir) al supermercado. \hfill
b) Tú \luecke{2.5cm} (estar) en el parque. \\[0.8em]
c) Ella \luecke{2.5cm} (ir) a la farmacia. \hfill
d) Nosotros \luecke{2.5cm} (estar) en la plaza. \\

\vspace{10pt}

% ===============================================================
% AUFGABE 3: Lückentext Wegbeschreibung
% ===============================================================

\begin{aufgabe}{3 -- Wegbeschreibung: Ergänze die Lücken}
Wähle aus: \es{todo recto}, \es{a la derecha}, \es{a la izquierda}, \es{cerca de}, \es{enfrente de}, \es{al lado de} \punktebox{6}
\end{aufgabe}

\noindent
a) -- ¿Dónde está el banco?\\
\hspace*{0.5cm} -- El banco está \luecke{4cm} la iglesia. (neben)

\vspace{0.8em}

\noindent
b) -- ¿Cómo llego al cine?\\
\hspace*{0.5cm} -- Sigue \luecke{4cm} y luego gira \luecke{4cm}. (geradeaus, nach rechts)

\vspace{0.8em}

\noindent
c) -- ¿Está lejos la farmacia?\\
\hspace*{0.5cm} -- No, está \luecke{4cm} el hospital. (in der Nähe von)

\vspace{0.8em}

\noindent
d) -- ¿Dónde está el museo?\\
\hspace*{0.5cm} -- Está \luecke{4cm} el parque. (gegenüber von)

\vspace{10pt}

% ===============================================================
% AUFGABE 4: Dialog schreiben
% ===============================================================

\begin{aufgabe}{4 -- Schreibe einen Dialog: Wegbeschreibung}
A fragt nach dem Weg zum Kino. B beschreibt den Weg. \punktebox{6}
\end{aufgabe}

\noindent
\textbf{A:} \luecke{10cm}

\vspace{0.8em}

\noindent
\textbf{B:} \luecke{10cm}

\vspace{0.8em}

\noindent
\textbf{A:} ¿Está lejos?

\vspace{0.8em}

\noindent
\textbf{B:} \luecke{10cm}

\vspace{0.8em}

\noindent
\textbf{A:} ¡Gracias!

\vspace{0.8em}

\noindent
\textbf{B:} \luecke{6cm}

% ===============================================================
% PUNKTEÜBERSICHT
% ===============================================================

\vspace{20pt}
\hrule
\vspace{10pt}

\begin{flushright}
\textbf{Punkte:} \luecke{1cm} / \maxpunkte
\end{flushright}

\end{document}
